%
% mittelwert.tex -- Illustration des Mittelwertsatzes
%
% (c) 2021 Prof Dr Andreas Müller, OST Ostschweizer Fachhochschule
%
\documentclass[tikz]{standalone}
\usepackage{amsmath}
\usepackage{times}
\usepackage{txfonts}
\usepackage{pgfplots}
\usepackage{csvsimple}
\usetikzlibrary{arrows,intersections,math,calc}
\definecolor{darkred}{rgb}{0.8,0,0}
\begin{document}
\def\skala{1}
\def\a{1}
\def\b{7}
\begin{tikzpicture}[>=latex,thick,scale=\skala,
declare function = {
f(\X) = 0.15*(\X-2)*(\X-6)*(\X-4)+0.3*\X+1;
}]

\pgfmathparse{atan(21/20)}
\xdef\w{\pgfmathresult}

\pgfmathparse{4-sqrt(3)}
\xdef\xone{\pgfmathresult}

\draw[color=blue] (\xone,{f(\xone)}) -- ++(\w:1.5);
\draw[color=blue] (\xone,{f(\xone)}) -- ++({180+\w}:1.5);

\pgfmathparse{4+sqrt(3)}
\xdef\xtwo{\pgfmathresult}

\draw[color=blue] (\xtwo,{f(\xtwo)}) -- ++(\w:1.5);
\draw[color=blue] (\xtwo,{f(\xtwo)}) -- ++({180+\w}:1.5);

\draw (\a,-0.05) -- ++(0,0.1);
\draw (\b,-0.05) -- ++(0,0.1);
\draw[line width=0.2pt] (\a,0) -- (\a,{f(\a)});
\draw[line width=0.2pt] (\b,0) -- (\b,{f(\b)});

\draw[color=darkred,line width=1.2pt] plot[domain=1:7,samples=100]
	({\x},{f(\x)});

\draw (\xone,-0.05) -- ++(0,0.1);
\draw[line width=0.2pt] (\xone,0) -- (\xone,{f(\xone)});
\node at (\xone,0) [below] {$\xi_1$};
\node[color=darkred] at ({\xone+0.1},{f(\xone)-0.1}) [above left] {$f(\xi_1)$};
\node[color=blue] at ({\xone-0.8},{f(\xone)-1.0}) [above left] {$f'(\xi_1)$};

\draw (\xtwo,-0.05) -- ++(0,0.1);
\draw[line width=0.2pt] (\xtwo,0) -- (\xtwo,{f(\xtwo)});
\node at (\xtwo,0) [below] {$\xi_2$};
\node[color=darkred] at ({\xtwo+0.1},{f(\xtwo)-0.1}) [above left] {$f(\xi_2)$};
\node[color=blue] at ({\xtwo-1.2},{f(\xtwo)-1}) [below right] {$f'(\xi_2)$};

\fill[color=white] (\xone,{f(\xone)}) circle[radius=0.05];
\draw[color=darkred] (\xone,{f(\xone)}) circle[radius=0.05];
\fill[color=white] (\xtwo,{f(\xtwo)}) circle[radius=0.05];
\draw[color=darkred] (\xtwo,{f(\xtwo)}) circle[radius=0.05];

\draw[color=blue] (\a,{f(\a)}) -- (\b,{f(\b)});

\fill[color=white] (\a,{f(\a)}) circle[radius=0.05];
\draw[color=darkred] (\a,{f(\a)}) circle[radius=0.05];
\node at (\a,0) [above] {$a$};
\node at (\a,{f(\a)}) [left] {$f(a)$};

\fill[color=white] (\b,{f(\b)}) circle[radius=0.05];
\draw[color=darkred] (\b,{f(\b)}) circle[radius=0.05];
\node at (\b,0) [below] {$b$};
\node at (\b,{f(\b)}) [right] {$f(b)$};

\draw[->] (-0.1,0) -- (8.3,0) coordinate[label={$x$}];
\draw[->] (0,-1) -- (0,5.7) coordinate[label={right:$y$}];

\node[color=blue] at ($0.2*(\a,{f(\a)})+0.8*(\b,{f(\b)})+(0.2,-0.2)$)
	[above left] {$\displaystyle\frac{f(b)-f(a)}{b-a}$};

\end{tikzpicture}
\end{document}

